% Self-contained source. Build with Tectonic, or run XeLaTeX twice.
\documentclass[11pt,a4paper]{article}
\usepackage{fontspec}
\usepackage[english]{babel}
\usepackage{amsmath,amssymb,amsthm,mathtools,booktabs,microtype}
\usepackage[a4paper,left=27mm,right=27mm,top=24mm,bottom=24mm,headheight=14pt]{geometry}
\usepackage{xcolor,fancyhdr}
\definecolor{ink}{HTML}{234452}
\usepackage[unicode,colorlinks=true,linkcolor=ink,citecolor=ink,urlcolor=ink,
 pdfauthor={Mikita Aharodni},pdftitle={Positive regular fibers of a catenary compartmental model with two outputs}]{hyperref}
\newtheorem{theorem}{Theorem}
\newtheorem{lemma}[theorem]{Lemma}
\newtheorem{proposition}[theorem]{Proposition}
\newcommand{\R}{\mathbb R}
\newcommand{\id}{\mathbf 1}
\newcommand{\diag}{\operatorname{diag}}
\newcommand{\Res}{\operatorname{Res}}
\newcommand{\dd}{\mathrm d}
\DeclareMathOperator{\supp}{supp}
\setlength{\emergencystretch}{2em}
\setlength{\parskip}{.18em}
\raggedbottom
\fancyhf{}\fancyhead[L]{\small Positive regular fibers of a catenary model}
\fancyhead[R]{\small M. Aharodni}\fancyfoot[C]{\thepage}
\renewcommand{\headrulewidth}{.2pt}\pagestyle{fancy}
\begin{document}
\thispagestyle{plain}
\begin{center}
{\LARGE\bfseries Positive regular fibers of a catenary\\[3pt] compartmental model with two outputs\par}
\vspace{.7em}
{\large Mikita Aharodni\par}
\vspace{.25em}
{\small 25 November 2026\par}
\end{center}
\begin{abstract}
Consider a bidirected path of \(n\ge4\) compartments with an inlet at the first
compartment, amount measurements at the third and last compartments, and a
single leak at a known site \(\ell\). We determine the maximum number of
positive regular rate vectors consistent with the same two output trajectories.
The maximum is one if the leak is at either measured compartment and two at
every other leak location. The upper bound follows by expressing the distal
gain as a concave minimum on an affine interval of endpoint spectral weights
of a Jacobi matrix. A five-compartment example and two extension operations
give sharpness for interior leaks at every chain length. For a single leak
at the penultimate compartment and \(n\ge5\), the two realizations have
different elimination rates on a nonempty open set of data. With the two
coprime spectral polynomials and the proximal gain fixed, at most one distal
gain admits two regular solutions with the same elimination rate. Such a
pair exists already for \(n=5\).
\end{abstract}
\noindent\textit{Keywords.} Structural identifiability, compartmental models,
inverse spectral problems, Jacobi matrices, positive fibers.

\section{Model and main result}\label{sec:model}
For \(1\le r<n\), write \(a_r>0\) for the transfer rate from \(r\) to \(r+1\)
and \(b_r>0\) for the reverse rate. The single leak has unknown rate \(h>0\)
at a prescribed site \(\ell\). With \(a_n=b_0=0\), the compartment matrix
has entries
\[
 A_{r+1,r}=a_r,\qquad A_{r,r+1}=b_r,\qquad
 A_{rr}=-a_r-b_{r-1}-h\id_{\{r=\ell\}}.
\]
All other entries vanish. The experiment is \(x'=Ax+e_1u,\ x(0)=0\), with
known continuous input \(u\) and observed amounts \(x_3,x_n\). Rates form the
vector \(\theta=(a_1,\ldots,a_{n-1},b_1,\ldots,b_{n-1},h)\).

Let \(D(s)=\det(sI-A)\), let \(W(s)\) be the characteristic polynomial of
the principal block on sites \(4,\ldots,n\), and put
\(p=a_1a_2,\ q=\prod_{r=1}^{n-1}a_r\). Cofactor expansion along the path
gives the two transfer functions
\begin{equation}\label{eq:transfer}
 e_3^\top(sI-A)^{-1}e_1=\frac{pW(s)}{D(s)},\qquad
 e_n^\top(sI-A)^{-1}e_1=\frac{q}{D(s)} .
\end{equation}
The second numerator is a positive constant. It therefore determines \(q\)
and the monic \(D\) from the reduced fraction. Multiplying the first transfer
function by \(D\) recovers \(pW\), whose leading coefficient is \(p\).
These conclusions include cancellations in the first transfer function.
Equality of the two transfer functions is equivalent to equality of the two
zero-state responses for every input. For the converse it even suffices to
use the constant unit input and differentiate its step response.

Define \(F_{n,\ell}(\theta)\) by listing the nonleading coefficients of
\(D,W\), each in \emph{ascending} order, followed by \(p,q\). Its domain
is \(\Theta=(0,\infty)^{2n-1}\) and its target has dimension
\(n+(n-3)+2=2n-1\). We call a point regular if
\(\det \dd F_{n,\ell}(\theta)\ne0\). Define
\[
 M(n,\ell)=\max_y\#\{\theta\in\Theta:F_{n,\ell}(\theta)=y,\
                           \det\dd F_{n,\ell}(\theta)\ne0\}.
\]
\begin{theorem}\label{thm:main}
For every \(n\ge4\) and prescribed leak location \(1\le\ell\le n\),
\[
 M(n,\ell)=
 \begin{cases}
  1,&\ell\in\{3,n\},\\
  2,&\ell\notin\{3,n\}.
 \end{cases}
\]
In the second case, two distinct positive regular preimages persist on a
nonempty open set of data.
\end{theorem}

Local identifiability on bidirected trees with a single input and a single
output is characterized in \cite[Theorem~5.2]{bortner}. Multiple observation
sites on paths are discussed in \cite[\S6]{ahmed}. A placement criterion for
local identifiability and its formal verification are given in the companion
note \cite{placement}. Here the inlet and outputs are fixed. The quantity in
Theorem~\ref{thm:main} counts regular points in the positive real domain and
takes a maximum over attainable data. The usual identifiability degree and the
general question in \cite[Question~5.3]{meshkat} concern a different count.

\section{Physical coordinates and minors}
Let \(c_r=\sqrt{a_rb_r}\), \(\beta_r=c_r^2\). The matrix \(-A\) is
diagonally similar to a symmetric tridiagonal matrix \(T\) with off-diagonal
entries \(-c_r\). It is positive definite since
\[
 v^\top T v=\sum_{r=1}^{n-1}
   (\sqrt{a_r}v_r-\sqrt{b_r}v_{r+1})^2+h v_\ell^2.
\]
If the right-hand side vanishes, the value \(v_\ell=0\) propagates along
every edge. Conversely, for any positive definite tridiagonal \(T\) with
strictly negative off-diagonals \(-c_r\), set \(w=T^{-1}e_\ell\).
Every entry of \(w\) is positive. Indeed, for sufficiently large \(z>0\),
\(I-T/z\) is entrywise nonnegative and has spectral radius less than one.
Its powers contain a positive contribution along every path, so the Neumann
series for \(T^{-1}\) is strictly positive. The inverse correspondence is
\begin{equation}\label{eq:gauge}
 a_r=c_r w_{r+1}/w_r,\qquad b_r=c_r w_r/w_{r+1},\qquad h=1/w_\ell .
\end{equation}
In fact \(A=-\diag(w)T\diag(w)^{-1}\) has column sums
\(-\id_{\{r=\ell\}}/w_\ell\). Conversely any positive scaling producing a
physical matrix solves \(Tw=t e_\ell\) for some \(t>0\), and scaling by
\(1/t\) gives the above normalization. The two coordinate changes are smooth.

Write \(L_j(s),R_j(s)\) for the determinants of \(sI+T\) on the intervals
\([1,j]\), \([j,n]\). Empty determinants are one. Put \(E=L_2,\ H=L_{n-1}\),
\(K=\prod_{r=3}^{n-1}\beta_r\), so \(W=R_4\). For \(\ell<n\), let \(C_\ell\)
denote the determinant on \([\ell+1,n-1]\). The following relations will be used:
\begin{align}
 p&=E(0) &&(\ell\ge3),\label{eq:p}\\
 K&=qW(0)/p &&(\ell\le3),\label{eq:K}\\
 q&=H(0)-D(0)C_\ell(0)/R_{\ell+1}(0) &&(\ell<n),\label{eq:q}\\
 q&=H(0) &&(\ell=n).\nonumber
\end{align}
For completeness, the three-term determinant recurrence at zero gives
\(L_{j-1}(0)=\prod_{r<j} a_r\) for \(j\le\ell\) and
\(R_{j+1}(0)=\prod_{r=j}^{n-1} b_r\) for \(j\ge\ell\). This proves
\eqref{eq:p}, \eqref{eq:K} and the terminal case. For \(\ell<n\) apply
the complementary-minor identity
\[
 H R_{\ell+1}-D C_\ell=L_{\ell-1}\prod_{r=\ell}^{n-1}\beta_r
\]
at zero and divide by \(R_{\ell+1}(0)>0\). The same identity, applied to sites
3 and \(n\), gives
\begin{equation}\label{eq:minor}
 WH-DB=K E,\qquad B=\det((sI+T)_{[4,n-1]}).
\end{equation}
Both identities also follow from Desnanot--Jacobi applied to the principal
interval between the distinguished sites after eliminating its exterior
blocks, or directly by the tridiagonal recurrence.

\section{The two-dimensional spectral fiber}
Assume throughout this section that \(\gcd(D,W)=1\). Write
\(0<\lambda_1<\cdots<\lambda_n\) for the eigenvalues of \(T\). Irreducible
tridiagonality makes each eigenvalue simple and each endpoint component of
an eigenvector nonzero, by the three-term recurrence. The spectral weights
at \(e_n\) satisfy
\[
 \rho_i>0,\qquad \sum_i\rho_i=1,\qquad
 H(s)/D(s)=\sum_i \rho_i/(s+\lambda_i).
\]
With \(P=K E\), identity \eqref{eq:minor} and residues give
\begin{equation}\label{eq:weights}
 \rho_i=\frac{P(-\lambda_i)}{W(-\lambda_i)D'(-\lambda_i)},\qquad
 \sum_i\frac{P(-\lambda_i)}{W(-\lambda_i)D'(-\lambda_i)}=1 .
\end{equation}
\begin{lemma}\label{lem:coords}
Let \(D,W\) be real monic coprime polynomials of degrees \(n,n-3\), with
\(D(s)=\prod_{i=1}^n(s+\lambda_i)\) and \(0<\lambda_1<\cdots<\lambda_n\).
A positive physical realization exists exactly when a polynomial \(P\) of
degree at most two satisfies \eqref{eq:weights} with every \(\rho_i>0\).
For any prescribed leak location, these polynomials are in bijection with
the physical realizations. When nonempty, they form an open convex subset
of an affine plane. Writing \(E(s)=s^2+es+f\), the coefficients
of \(D,W\) together with \(e,f\) are local smooth coordinates.
\end{lemma}
\begin{proof}
Given such weights, orthogonalize \(1,x,\ldots,x^{n-1}\) for the discrete
measure \(\sum_i\rho_i\delta_{\lambda_i}\). The multiplication operator by
\(x\) gives the unique Jacobi matrix with the specified endpoint measure,
after reversing indices and alternating basis signs to obtain the convention
for \(T\). Its eigenvalues are positive. For \(0\le k\le n-4\),
\[
 \sum_i\rho_i W(-\lambda_i)(-\lambda_i)^k
   =\sum_i\frac{P(-\lambda_i)(-\lambda_i)^k}{D'(-\lambda_i)}=0
\]
by the leading coefficient of Lagrange interpolation, since the numerator
has degree at most \(n-2\). Thus the monic orthogonal polynomial of degree
\(n-3\) is \((-1)^{n-3}W(-x)\). It is the characteristic polynomial of the
first \(n-3\) sites in reversed order, proving that the reconstructed tail
has polynomial \(W(s)\). Its own polynomial \(K E\) in \eqref{eq:minor}
agrees with \(P\) at \(n\) distinct nodes and hence identically. In
particular the coefficient of \(s^2\) in \(P\) is strictly positive.

With \(E\) monic, put
\begin{equation}\label{eq:alpha}
 \alpha_i=\frac{E(-\lambda_i)}{W(-\lambda_i)D'(-\lambda_i)},\qquad
 S=\sum_i\alpha_i,\qquad \rho_i=\alpha_i/S,\quad K=1/S.
\end{equation}
The sorted simple roots, \(\alpha_i\) and the coefficients of the
orthonormal polynomials depend smoothly on nearby data. The Gram determinants
are positive and the signs of the off-diagonals are prescribed. Positivity
of all \(\alpha_i\) persists under small independent changes of \(D,W,e,f\),
and the construction just given yields a smooth inverse via \eqref{eq:gauge}.
\end{proof}
The inverse spectral problem with the two spectra prescribed is also treated
by Castillo and Gordillo-Núñez \cite[Lemmas~2.1, 3.1 and Theorem~3.1]{castillo}.
The remaining gain conditions determine the count below.

\section{A variational upper bound}
For \(0\le m\le n-1\), set
\begin{equation}\label{eq:phi}
 \Phi_m(\rho)=\min_{\deg r\le m,\ r(0)=1}
   \sum_{i=1}^n\frac{\rho_i}{\lambda_i}\,r(\lambda_i)^2 .
\end{equation}
\begin{lemma}\label{lem:phi}
The distal gain satisfies \(q/D(0)=\Phi_{n-\ell}(\rho)\).
For fixed nodes, \(\Phi_m\) is concave in the weights and real analytic on
the positive orthant.
\end{lemma}
\begin{proof}
Write \(r(x)=1-xv(x)\) with \(\deg v<m\) and \(z=v(T)e_n\). Its energy is
\(e_n^\top T^{-1}e_n-2e_n^\top z+z^\top Tz\). The Krylov columns
\(e_n,\ldots,T^{m-1}e_n\) span the last \(m\) coordinates: column \(j\),
starting at \(j=0\), first reaches site \(n-j\) with nonzero coefficient
\((-1)^j\prod_{t=n-j}^{n-1}c_t\). The minimum subtracts from
\(e_n^\top T^{-1}e_n\) the last diagonal entry of the inverse of the trailing
principal block of order \(m\), taking this correction as zero when \(m=0\).
With \(m=n-\ell\), the two entries are \(H(0)/D(0)\) and
\(C_\ell(0)/R_{\ell+1}(0)\). Equation~\eqref{eq:q} and its terminal case give the claim.

Each candidate energy is linear in \(\rho\), so its infimum is concave.
For \(\mu_j=\sum_i\rho_i\lambda_i^j\), \(b_j=\mu_j\), and
\(G_{jk}=\mu_{j+k+1}\) (\(0\le j,k<m\)), completing the square gives
\begin{equation}\label{eq:gram}
 \Phi_m=\mu_{-1}-b^\top G^{-1}b .
\end{equation}
The weighted Vandermonde Gram matrix \(G\) is positive definite for distinct
positive nodes and positive weights. This also proves attainment and analyticity.
\end{proof}

\begin{proposition}\label{prop:upper}
The number of positive regular preimages is at most two at every data point,
and at most one if \(\ell=3\) or \(\ell=n\).
\end{proposition}
\begin{proof}
First suppose that \(D,W\) are coprime. If \(\ell\ge3\), fixing \(p\) imposes
\(P(0)=p[s^2]P\) by \eqref{eq:p}. This homogeneous nonzero linear equation
is independent of the normalization in \eqref{eq:weights} whenever they have
a common solution. Their positive intersection is an affine interval.
If \(\ell\le2\), fixing \(p,q\) fixes \(K=qW(0)/p=[s^2]P\). This also cuts
out an affine interval unless it is constant on the normalization plane.
In the latter case the pair \((K,q)\), equivalently \((p,q)\), has rank at
most one on that plane. Such points are singular by Lemma~\ref{lem:coords}.

On the resulting interval of weights the remaining equation is
\(\Phi_{n-\ell}=q/D(0)\). Three distinct points at this level give zero
secant slopes on either side of the middle point. Concavity and
differentiability force its derivative to vanish there, contrary to
regularity in the coordinates of Lemma~\ref{lem:coords}. When \(\ell=3\),
both \(P(0)=p[s^2]P\) and \([s^2]P=qW(0)/p\) are linear equations.
When \(\ell=n\), the second observation is \(q=H(0)\), also linear in \(P\)
by \eqref{eq:weights}. In either case an affine fiber through two positive
points contains a segment, so neither endpoint is locally isolated.

Finally, choose any finite collection of distinct positive regular
preimages at the same data. By the inverse function theorem they continue
in disjoint positive neighborhoods over one common open set of data.
The resultant \(\Res(D,W)\) is a nonzero polynomial on the ambient coefficient
space, so this open set contains coprime data. The bounds already proved
apply to the collection.
\end{proof}

\section{Leaks at sites \texorpdfstring{\(1,2,3,n\)}{1, 2, 3 and n}}\label{sec:boundary}
Set \(T_0=n,\ t_j=j-(n+1)/2,\ \lambda_j=T_0+t_j\) for \(1\le j\le n\),
and let \(a_*=(n-3)/2,\ b_*=(n-1)/2\). Put
\(\mu_j=T_0+j-n/2\) for \(2\le j\le n-2\). These are the midpoints of
the gaps between the \(\lambda_j\) other than the two outermost gaps. Use
positive spectral variable \(x\) and write
\[
 \widehat D(x)=\prod_{j=1}^n(x-\lambda_j),\qquad
 \widehat W(x)=\prod_{j=2}^{n-2}(x-\mu_j),\qquad
 \Delta_j=\widehat W(\lambda_j)\widehat D'(\lambda_j).
\]
Take \(\widehat E_u(x)=(x-T_0)^2+u(x-T_0)-a_*b_*\), \(-1<u<1\).
In the \(s\)-variable \(E_u(s)=\widehat E_u(-s)\).
The \(\Delta_j\) are positive at the two endpoints and negative at every
interior node, by counting factors to the right. The same signs hold for
\(\widehat E_u(\lambda_j)\): at \(u=\pm1\) its roots in \(t=x-T_0\)
are \((-b_*,a_*)\) and \((-a_*,b_*)\), respectively. The intermediate signs follow by interpolation in \(u\).
Thus \(\alpha_j=\widehat E_u(\lambda_j)/\Delta_j>0\). Reflection gives
\(\Delta_{n+1-j}=\Delta_j\) and \(\sum_j t_j/\Delta_j=0\), so \(S\) and
\(K=1/S\) are constant and the normalized weights are affine in \(u\).

For \(\ell=1,2\), \(m=n-\ell\ge n-2\). Exactly \(n-2\) support nodes survive
at either end of the interval. The polynomial
\(\prod_{\lambda\in \supp}(1-x/\lambda)\) on that limiting support is an
admissible test polynomial with limiting energy zero. In the interior the
attained minimum is positive, since no polynomial of degree \(m<n\)
normalized at zero vanishes at all \(n\) nodes. Concavity now gives two
crossings with nonzero derivative at a level between zero and the maximum.
They have the same \(q\) and \(K\), hence the same \(p=qW(0)/K\).

To verify the full Jacobian, allow
\(\widehat E_{u,v}(x)=(x-T_0)^2+u(x-T_0)+v\).
Then \(S=c_2+c_0v\), and the residues of
\(1/(\widehat D\,\widehat W)\) give
\[
 c_0=\sum_j\frac1{\Delta_j}
 =-\sum_{j=2}^{n-2}\frac1{\widehat D(\mu_j)\widehat W'(\mu_j)}<0.
\]
Both factors in each product on the right have the same sign. Consequently
\(K_v=-c_0/S^2\ne0\), while \(K_u=0\) and \(q_u\ne0\) at the crossings.
Lemma~\ref{lem:coords} proves their full regularity.

For \(\ell=3,n\), use \(u=0,\ v=-a_*b_*\). Now
\(p=T_0^2-T_0u+v\). At \(\ell=3\) the independent differentials of \(p,K\)
give a regular point through \(q=Kp/W(0)\). At \(\ell=n\), vary both
variables with \(\dd v/\dd u=T_0\) to hold \(p\) fixed. Then
\[
 \frac{\dd\alpha_j}{\dd u}=\frac{\lambda_j}{\Delta_j},\qquad
 \frac{\dd \Phi_0}{\dd u}=\frac{c_0}{S}(1-T_0\Phi_0).
\]
Symmetry of the normalized weights gives
\(T_0\Phi_0=\sum_j\rho_j T_0^2/(T_0^2-t_j^2)>1\): each term in the weighted
average is at least one and the endpoint weights are positive. This supplies
a regular point for the other terminal case.

\section{Interior leaks and a five-compartment seed}
Suppose \(4\le\ell<n\). Hold \(D,W,p\) fixed and set \(E_t(s)=s^2+ts+p\).
Use \eqref{eq:alpha} to define \(\alpha_i(t),S(t),\rho_i(t)\), and abbreviate
\(\Phi_m(\rho(t))\) to \(\phi(t)\), \(m=n-\ell\).
\begin{lemma}\label{lem:peak}
If \(\alpha_i(t)>0\) on \([a,c]\) for every \(i\), and \(a<b<c\) satisfy
\(\phi(b)>\max\{\phi(a),\phi(c)\}\), then there exist two distinct positive
regular preimages of a common data vector.
\end{lemma}
\begin{proof}
The map \(t\mapsto\rho(t)\) has the form \((A+tB)/(C+tD_1)\). It is injective:
identical weights give identical \(P=KE_t\) by polynomial interpolation at
the \(n\) nodes, and normalization of \(E_t\) then recovers \(t\).
Its image is an affine interval and its speed is nonzero since the numerator
of its derivative is the constant vector \(BC-AD_1\). In an affine
coordinate along that interval \(\Phi_m\) is concave. Any level strictly
between \(\max\{\phi(a),\phi(c)\}\) and \(\phi(b)\) has a crossing on each
side of the middle point. The secant slope to the middle point bounds the
derivative away from zero at each crossing. Since \(p=E_t(0)\) is the other
independent coordinate at fixed \(D,W\), nonvanishing of \(\dd q/\dd t\)
is precisely the remaining full-rank condition. Apply Lemma~\ref{lem:coords}.
\end{proof}

For \((n,\ell)=(5,4)\), take
\begin{align*}
 D(s)&=s^5+17s^4+96s^3+206s^2+150s+24,\\
 W(s)&=s^2+7s+8,\qquad p=2.
\end{align*}
These polynomials are coprime, as \(D\bmod W=6s+24\) and \(W(-4)=-4\).
For \(8\le k\le12\), where \(k\) is a scalar parameter, define
\begin{align*}
 P_k&=ks^2+(11k/2+6)s+2k,\\
 H_k&=s^4+(k/4+12)s^3+(5k/2+38)s^2+(9k/2+36)s+k+6.
\end{align*}
Direct multiplication gives
\(WH_k-D(s+2+k/4)=P_k\). The two positive physical rate vectors
\begin{align}
 \theta_A&=(1,2,3,1,\ 3,1,2,2,\ 2),\nonumber\\
 \theta_B&=\left(\tfrac{16}{29},\tfrac{29}{8},\tfrac94,\tfrac43,\
                 \tfrac{481}{232},\tfrac32,\tfrac89,3,\ \tfrac{16}{9}\right)
                 \label{eq:seed}
\end{align}
give these \(D,W,p\), \(q=6\) and \(k=12,8\), respectively. This can be
checked by the tridiagonal determinant recurrence. Their spectral weights
are positive, and \eqref{eq:weights} is affine in \(k\), so the weights
remain normalized and positive on the entire segment. The last diagonal
entry of \(T\) is the difference of the respective subleading coefficients
of \(D,H_k\), namely \(5-k/4\). Formula \eqref{eq:q} gives
\[
 q(k)=k+6-\frac{24}{5-k/4}.
\]
With \(t=11/2+6/k\), the three points \(a=6,\ b=61/10,\ c=25/4\)
correspond to \(k=12,10,8\). Division by \(D(0)=24\) yields
\[
 (\phi(a),\phi(b),\phi(c))=(1/4,\,4/15,\,1/4).
\]
This proves the strict inequality in Lemma~\ref{lem:peak}. In the coordinate
order of Section~\ref{sec:model} the full Jacobian determinants at \(\theta_A,\theta_B\)
are respectively \(3456\) and \(-4698\).

\section{Two extension operations}
Start with an interior placement satisfying Lemma~\ref{lem:peak}. Adjoin a
large node \(L\) and a new tail eigenvalue \(L-\delta>\max_i\lambda_i\), setting
\[
 D^+(s)=(s+L)D(s),\qquad W^+(s)=(s+L-\delta)W(s).
\]
Keep \(E_t\) unchanged. Formula \eqref{eq:alpha} gives exactly
\begin{equation}\label{eq:extend}
 \alpha_i^+(t)=\frac{\alpha_i(t)}
 {(L-\lambda_i)(L-\delta-\lambda_i)},\qquad
 \alpha_L^+(t)=\frac{E_t(-L)}{-\delta W(-L)D(-L)} .
\end{equation}
For sufficiently large \(L\) all weights are positive throughout \([a,c]\).
Indeed, \(E_t(-L)\) is uniformly positive there, and the sign of
\(W(-L)D(-L)\) is negative by their degrees. Coprimality persists for the
choices below, since the new node and tail eigenvalue are distinct and exceed all
old nodes and tail eigenvalues. The latter lie in the spectral interval by
the principal submatrix interlacing inequalities. Lemma~\ref{lem:coords}
thus gives a positive physical chain for each \(t\).

\paragraph{Operation A: \((n,\ell)\mapsto(n+1,\ell+1)\).}
Set \(\delta=L/2\). At each of the three fixed values of \(t\) the old
normalized weights tend to their original values. The new normalized
weight has order \(O(L^{6-2n})\), by \eqref{eq:extend}. Here \(m=n-\ell\)
stays fixed. The contribution of the new node to the highest moment
in \eqref{eq:gram}, of order \(2m-1\), is \(O(L^{5-2\ell})=o(1)\).
All lower moments, including the inverse moment, converge. Invertibility
of the limiting Gram matrix gives convergence of the three minima, and the
strict inequality persists for a common finite \(L\).

\paragraph{Operation B: \((n,\ell)\mapsto(n+1,\ell)\).}
Fix \(\varepsilon>0\) and take \(\delta=L^{7-2n}/\varepsilon\).
At each fixed \(t\), \eqref{eq:extend} implies
\(L^2\alpha_i^+\to\alpha_i\) and \(L^2\alpha_L^+\to\varepsilon\). The new
node thus has limiting mass \(\varepsilon/(S+\varepsilon)>0\). The residual
degree increases from \(m\) to \(m+1\). We claim that
\begin{equation}\label{eq:limit}
 \Phi_{m+1}(\rho^+(t))\ \longrightarrow\
 \frac{S(t)}{S(t)+\varepsilon}\Phi_m(\rho(t)).
\end{equation}
For the upper limit, test with \((1-x/L)r_*(x)\), where \(r_*\) is an old
minimizer. The new-node energy vanishes and the old-node energy tends to
the right-hand side of \eqref{eq:limit}.

For the lower limit, choose a sequence realizing the lower limit of the
energies and minimizers \(r_L\) along this sequence. Their energies are
bounded by the preceding test. Every old weight has a positive limit, so the
values of \(r_L\) at all old nodes are bounded. Since \(m+1<n\), evaluation
at \(m+2\) fixed old nodes gives an invertible Vandermonde system, bounding
all coefficients. At the new node \(r_L(L)=O(\sqrt L)\) by its positive
limiting weight. After division by \(L^{m+1}\) the leading coefficient
therefore tends to zero. A subsequence converges coefficientwise to a
polynomial of degree at most \(m\) taking value one at zero. Its contribution
from the old nodes is at least the right-hand side of \eqref{eq:limit},
and the new-node contribution is nonnegative. This proves the lower bound.
Now choose one sufficiently small \(\varepsilon>0\) so that the three
limiting values retain the strict inequality, since each factor
\(S(t)/(S(t)+\varepsilon)\) tends to one as \(\varepsilon\downarrow0\).
Then choose a common sufficiently large finite \(L\).

\begin{proof}[Completion of the proof of Theorem~\ref{thm:main}]
Starting at \((5,4)\), apply A exactly \(\ell-4\) times to reach
\((\ell+1,\ell)\), then B exactly \(n-\ell-1\) times. This covers every
\(4\le\ell<n\). Lemma~\ref{lem:peak} gives sharpness there, and Section~\ref{sec:boundary}
supplies the other leak locations. Proposition~\ref{prop:upper} gives the
matching upper bounds. At any two distinct regular preimages the inverse
function theorem supplies two disjoint positive branches over a common
open neighborhood of their data.
\end{proof}

\section{The elimination rate at the penultimate compartment}\label{sec:preterminal}
An individual rate may take the same value in two different realizations.
Clemens et al.\ studied this question for mammillary models with one input,
one output and no leak \cite[Theorem~3.2]{clemens}. Huneke and Keen
\cite[Proposition~4.5]{huneke} proved that the leak rate is globally identifiable
in a strongly connected model with one input when the sole output is at the
sole leak site. Here the leak site lies strictly between the two outputs.

\begin{theorem}\label{thm:preterminal}
Let \(n\ge5\) and \(\ell=n-1\). On a nonempty Euclidean open set of coefficient
data \((D,W,p,q)\), two positive regular realizations have distinct elimination
rates \(h\). Their two output responses agree for every common input.

Fix \(D,W,p\) with \(\gcd(D,W)=1\). Whenever two positive regular solutions
occur at the same \(q\), their backward rates \(b_{n-1}\) are distinct.
At most one value of \(q\) on this slice admits two such solutions with
the same \(h\). This exception occurs for \(n=5\).
\end{theorem}

\begin{proof}
Put \(g=q/D(0)>0\). Let \(z,v,r\) denote the last diagonal entries of
\(T,T^2,T^{-1}\), respectively. They are the moments of orders \(1,2,-1\)
of the endpoint spectral measure. The last compartment has no leak, so
\(z=b_{n-1}\) and \(v=z^2+a_{n-1}b_{n-1}\).
Equation~\eqref{eq:q} at \(\ell=n-1\) gives \(g=r-1/z\).
For a constant unit input, the equilibrium \(x^*=(-A)^{-1}e_1\) satisfies
\(h x^*_{n-1}=1\) by mass balance. The last compartment satisfies
\(b_{n-1}x^*_n=a_{n-1}x^*_{n-1}\), with \(x^*_n=g\). It follows that
\begin{equation}\label{eq:leak}
 g=r-\frac1z,\qquad h=\frac{a_{n-1}}{g b_{n-1}}
                    =\frac{v-z^2}{g z^2}.
\end{equation}
Here \(h\) denotes a first-order rate constant. The instantaneous amount lost
per unit time is \(h x_{n-1}(t)\).

On a nonempty slice with coprime \(D,W\) and fixed \(p\), normalization in
\eqref{eq:weights} and \(P(0)=p[s^2]P\) cut out an affine line. The two
conditions are independent since the latter is homogeneous and the former
takes the value one. By Lemma~\ref{lem:coords}, positive realizations
correspond bijectively to an open interval on this line. Each of \(z,v,r\)
is affine in the spectral weights. If \(z\) is constant on the interval,
then \(g\) is affine there. Two solutions at one level would force \(g\) to
be constant on the entire interval, contradicting regularity. Thus on a slice
containing a regular pair, \(z\) is an affine coordinate and distinguishes
the two solutions. Write \(r=\alpha z+\beta\) and \(v=\gamma z+\delta\).
Since \(v>0\), the pair \((\gamma,\delta)\) is nonzero.

For distinct values \(x,y\) of \(z\) giving the same \(g\),
\eqref{eq:leak} gives \(\alpha=-1/(xy)\) and
\(1/x+1/y=\beta-g\). Subtracting the two expressions for \(h\) yields
\[
 g\bigl(h(x)-h(y)\bigr)
 =\left(\frac1x-\frac1y\right)\bigl(\gamma-\delta(g-\beta)\bigr).
\]
Hence the two elimination rates agree precisely when
\begin{equation}\label{eq:exception}
 \gamma=\delta(g-\beta).
\end{equation}
If \(\delta=0\), then \(\gamma\ne0\), and equality is impossible. Otherwise
the only candidate is \(q=D(0)(\beta+\gamma/\delta)\).

For each \(n\ge5\), Theorem~\ref{thm:main} supplies two regular inverse
branches over a common open set of data. For this leak location the
construction uses the pair \eqref{eq:seed} and \(n-5\) applications of
operation A. The open set contains data with coprime \(D,W\), since
\(\Res(D,W)\) is a nonzero polynomial in their coefficients. Fix such
\(D,W,p\) and vary \(q\) in the open set. At any level other than the one
allowed by \eqref{eq:exception}, the two values of \(h\) differ.
Continuity preserves this strict inequality in an open neighborhood
in the full coefficient space.

For an exceptional level take \(n=5,\ \ell=4\) and
\begin{align*}
 \theta_1&=(3,4,4,4,\ 3,2,4,1,\ 2),\\
 \theta_2&=\left(\tfrac{44}{17},\tfrac{51}{11},\tfrac{10}{3},\tfrac{24}{5},\
                 \tfrac{332}{187},\tfrac{11}{3},3,\tfrac65,\ 2\right).
\end{align*}
The determinant recurrence gives both vectors the data
\[
 D(s)=s^5+27s^4+230s^3+672s^2+560s+96,\quad
 W(s)=s^2+11s+6,\quad (p,q)=(12,192).
\]
Here \(D\bmod W=-64(4s+3)\) and \(W(-3/4)=-27/16\).
The two Jacobian determinants in the coordinate order of Section~\ref{sec:model}
are \(-2359296\) and \(3760128\). Both points have \(h=2\), while
\(b_4=1\) at the first and \(b_4=6/5\) at the second.
\end{proof}

Thus on a coprime slice with two regular solutions an independent exact
measurement of \(b_{n-1}\) distinguishes them. A measurement of \(h\)
distinguishes them at every level except a possible level \eqref{eq:exception}.

\section*{Computations and formal verification}
The \href{https://github.com/qqhgdkktu/catenary-identifiability/tree/global-multiplicity-v2/global-multiplicity}{source and verification package}
includes two exact rational calculations of the examples in Section~\ref{sec:preterminal}.
One uses matrix characteristic polynomials. The other differentiates the
tridiagonal determinant recurrence symbolically. Lean~4 verifies
\(M(4,2)=2\), \(M(5,4)\ge2\), several spectral identities for arbitrary
chain length, and the algebraic equivalence and uniqueness in
\eqref{eq:exception}. The existence of the two regular branches for every
length is proved analytically in Sections~5--7.

\begin{thebibliography}{9}
\bibitem{bortner} C. Bortner, E. Gross, N. Meshkat, A. Shiu and S. Sullivant.
Identifiability of linear compartmental tree models and a general formula
for input-output equations. \textit{Advances in Applied Mathematics}
\textbf{146} (2023), 102490.
\href{https://arxiv.org/abs/2106.08487v2}{arXiv:2106.08487v2}.
\bibitem{ahmed} S. Ahmed, N. Crepeau, P. R. Dessauer Jr., A. Edozie,
O. Garcia-Lopez, T. Grimsley, J. Lopez Garcia, V. Neri and A. Shiu.
Identifiability of directed-cycle and catenary linear compartmental models.
\textit{SIAM Journal on Applied Dynamical Systems} \textbf{25} (2026),
no.~1, 304--350.
\href{https://doi.org/10.1137/25M1728636}{doi:10.1137/25M1728636}.
\href{https://arxiv.org/abs/2412.05283v3}{arXiv:2412.05283v3}.
\bibitem{meshkat} N. Meshkat and A. Shiu. Structural identifiability of
compartmental models: recent progress and future directions.
\href{https://arxiv.org/abs/2507.04496v1}{arXiv:2507.04496v1} (2025).
\bibitem{castillo} K. Castillo and G. Gordillo-Núñez. Wendroff's theorem beyond
consecutive degrees and related inverse spectral problems.
\href{https://arxiv.org/abs/2602.20894v1}{arXiv:2602.20894v1} (2026).
\bibitem{clemens} K. Clemens, J. Martinez, A. Shiu, M. Thompson and B. Warren.
Parameter identifiability of linear-compartmental mammillary models.
\href{https://arxiv.org/abs/2506.21889v2}{arXiv:2506.21889v2} (2025).
\bibitem{huneke} T. Huneke and T. Keen. On the identifiability of linear
compartmental model parameters. Texas A\&M University, REU research report,
17 July 2025. \href{https://artsci.tamu.edu/mathematics/_files/_docs/reu/results/2025/huneke-keen-report.pdf}{Proposition~4.5}.
\bibitem{placement} M. Aharodni. Two outputs in a
catenary compartmental model. Research note and Lean formalization,
\href{https://github.com/qqhgdkktu/catenary-identifiability/tree/v0.1.0}{version 0.1.0} (2026).
\end{thebibliography}
\end{document}
